%% 第22回編集委員会（物理第2巻）：事務局。第1巻の ph1-front.tex と同じ部品・同じ形・同じページ数（各2ページ）
%% ph2-front.tex  前付の見開き2つ（物理シリーズ第2巻）
%%   地図は物理シリーズ全体のもの（第1巻と同じ表）。第1巻で力学の例を書いた所（相似の図、式の読み方の見本など）を、熱学の例にした。
%%   p.4〜5  保存則と相似の地図（第8回の構成案「保存則とバランス式の表」と「相似の表」を図にしたもの）  \phConservationMapPages
%%   p.6〜7  この本の使い方（数学 bk-howto.tex の物理版。物理で足した囲みの見本をのせる）            \phHowToPages
%% 章のページ（\phchappage{熱-2}）は ph2-pages.tex（gen_phys.py が作る）から読む。
\makeatletter
% 前付の見出し（数学の「この本の路線図」「この本の使い方」と同じ形）
\newcommand\ph@fronthead[2]{%
  \noindent{\sffamily\bfseries\fontsize{24}{30}\selectfont\color{bkDeep}#1}\hfill{\sffamily\fontsize{8.5}{11}\selectfont\color{bkGray}#2}\par\vspace{1mm}
  {\color{bkVol}\hrule height 1.4pt}\vspace{3.5mm}}
% 章の札（ページつき）。この巻（第2巻）の章はページも出す。ほかの巻の章は札だけ
\newcommand\ph@cp[1]{\ph@routestop{#1}}
\newcommand\ph@cps[1]{\foreach \c [count=\k] in {#1} {\ifnum\k>1 \hspace{0.25em}\fi\ph@cp{\c}}}

%% ============================================================
%% p.4〜5  保存則と相似の地図
%% ============================================================
\newcommand\phConservationMapPages{%
\clearpage\setcounter{page}{4}%
\global\bk@decofalse\bkfolio{true}\bksetrunning{}{}{}%
{\def\R##1##2{##1}\pdfbookmark[0]{保存則と相似の地図}{ph-front-map}\typeout{BKPAGE front 保存則と相似の地図 \thepage}}%
\ph@fronthead{保存則と相似の地図}{全5巻の背骨}
{\fontsize{10.5}{19}\selectfont
物理のたくさんの法則の多くは、たった1つの形の式から生まれます。それが下の「バランス式」で、おこづかい帳のお金の出し入れと同じ形です。\par}
\vspace{1mm}
\begin{bkfigure}
\begin{tikzpicture}[x=1mm, y=1mm]
  % 見えない境界（検査体積）
  \draw[bkTheme, line width=1.1pt, dash pattern=on 3pt off 2pt, rounded corners=3mm, fill=bkTheme!5] (38,0) rectangle (102,27);
  \node[anchor=north west, text=bkDeep, font=\sffamily\bfseries\footnotesize] at (40,26) {見えない境界で囲んだ所};
  % たまる量（水位）
  \fill[bkBlue!22] (40,2) rectangle (100,13);
  \draw[bkBlue, line width=1pt] (40,13) -- (100,13);
  \node[text=bkBlue!70!black, font=\sffamily\bfseries\small] at (70,7.5) {たまる量（増える・減る）};
  % 生じる量
  \node[draw=phFlowR, line width=0.9pt, fill=white, star, star points=8, star point ratio=1.6, minimum size=6mm, inner sep=0pt] (g) at (55,18) {};
  \node[anchor=west, text=phFlowR!80!black, font=\sffamily\bfseries\small] at (60,18) {生じる量};
  % 入る・出る
  \draw[-{Stealth[length=3.6mm, width=3mm]}, line width=3.2pt, bkTheme] (6,15) -- (37,15);
  \node[anchor=south, text=bkDeep, font=\sffamily\bfseries\small] at (20,17) {入る量};
  \draw[-{Stealth[length=3.6mm, width=3mm]}, line width=3.2pt, bkGray] (103,15) -- (134,15);
  \node[anchor=south, text=bkInk, font=\sffamily\bfseries\small] at (119,17) {出る量};
\end{tikzpicture}
\end{bkfigure}
\vspace{1mm}
\begin{center}
\tikz\node[fill=bkTheme, text=white, rounded corners=2mm, inner xsep=5mm, inner ysep=2.6mm, font=\sffamily\bfseries\large]
  {入る量 \,$-$\, 出る量 \,$+$\, 生じる量 \,$=$\, たまる量};
\end{center}
\vspace{1mm}
{\fontsize{10.5}{19}\selectfont
「数える量」をかえると、力学・流体・熱・電気・原子炉の法則になります。生じる量が0の量は「保存する量」です。札の横の数は、この巻のページです。\par}
\vspace{2.5mm}
{\footnotesize\noindent
\begin{tblr}{colspec={Q[l,m,wd=19mm] Q[l,m,wd=25mm] Q[l,m,wd=17mm] X[l,m] Q[l,m,wd=36mm]}, width=\linewidth,
  row{1}={bg=bkVol!14, font=\sffamily\bfseries\footnotesize}, hlines={0.4pt,bkLine}, rowsep=0.5pt, colsep=2.5pt,
  column{1}={font=\sffamily\bfseries}}
数える量 & 入る・出るもの & 生じる量 & 法則・式 & 章 \\
質量 & 質量流量 & 0 & 連続の式 & \ph@cps{熱-4,力-13} \\
運動量 & 力・運動量の流れ & 0 & 運動方程式、運動量保存、運動量の法則 & {\ph@cps{力-2,力-3}\\ \ph@cps{力-13}} \\
角運動量 & トルク & 0 & 回転の運動方程式、ポンプのオイラーの式 & \ph@cps{力-6,力-13} \\
エネルギー & 熱・仕事・エンタルピーの流れ・電磁場の流れ & 0 & 第一法則、定常流動エネルギー式、熱伝導のバランス、力学的エネルギー保存、ベルヌーイの式、ポインティングベクトル & {\ph@cps{熱-2,熱-4,熱-11}\\ \ph@cps{力-4,力-13,電-14}} \\
電荷 & 電流 & 0 & 電荷保存、キルヒホッフの電流則 & \ph@cps{電-4} \\
水分 & 水蒸気の流れ・結露・加湿 & 0 & 水分のバランス（空気線図） & \ph@cps{熱-9} \\
エントロピー & 熱÷温度 & 0以上 & 第二法則 & \ph@cps{熱-6,熱-15} \\
中性子 & 中性子の流れ & 生まれる数－吸収される数 & 中性子のバランス、増倍率 & \ph@cps{原-15} \\
\end{tblr}\par}
\clearpage
%% ---------------- p.5 相似 ----------------
\ph@fronthead{相似（にた形の式）の地図}{分野が違っても、式は同じ形}
{\fontsize{10.5}{19}\selectfont
お湯の冷め方と、コンデンサーに電気がたまる回路は、同じ形の式で書けます（一方を解けば、もう一方も解ける）。分野が違っても式が同じ形になることを、\textbf{相似}とよびます。\par}
\vspace{1mm}
\begin{bkfigure}
\begin{tikzpicture}[x=1mm, y=1mm]
  \path[use as bounding box] (0,-17.5) rectangle (139,32);
  % 左：ヒーターで温めるお湯の入れ物（熱容量）と、周りへ熱が逃げる壁（熱抵抗）
  \fill[bkLine!35, rounded corners=1.5mm] (0,1) rectangle (58,31);
  \node[anchor=north west, font=\sffamily\footnotesize, text=bkInk] at (1,30.4) {周り $T_{\mathrm a}$};
  \draw[bkInk, line width=2.6pt, line cap=round] (10,24) -- (10,4) -- (36,4) -- (36,24);
  \fill[bkBlue!22] (11.3,5.3) rectangle (34.7,19);
  \draw[bkBlue, line width=0.9pt] (11.3,19) -- (34.7,19);
  \node[font=\sffamily\footnotesize, text=bkBlue!70!black] at (23,15.8) {お湯 $T$};
  \node[font=\sffamily\footnotesize, text=bkBlue!70!black] at (23,11.8) {熱容量 $C_{\mathrm{th}}$};
  \draw[phFlowR, line width=0.9pt, decorate, decoration={coil, aspect=0.5, segment length=1.7mm, amplitude=1.1mm}] (15,7.4) -- (31,7.4);
  \node[anchor=north, font=\sffamily\footnotesize] at (23,3.4) {ヒーター $P$};
  \draw[-{Stealth[length=2.6mm]}, line width=1.3pt, phFlowR, decorate, decoration={snake, amplitude=0.7mm, segment length=3mm, post length=2.6mm}] (37.5,12) -- (52,12);
  \node[anchor=south, font=\sffamily\footnotesize, inner sep=1pt] at (45,14) {逃げる熱};
  \node[anchor=north east, font=\sffamily\footnotesize, inner sep=1pt] at (57.6,30.6) {壁の熱抵抗 $R_{\mathrm{th}}$};
  \draw[bkGray, line width=0.5pt] (44,26.4) -- (36.8,21);
  \node[anchor=north, align=center, font=\small] at (29,-1) {$C_{\mathrm{th}}\dfrac{\dd T}{\dd t}+\dfrac{T-T_{\mathrm a}}{R_{\mathrm{th}}}=P$\\[1mm]{\footnotesize $T-T_{\mathrm a}$（周りとの温度差）}};
  % まん中の等号
  \node[font=\sffamily\bfseries\Large, text=bkDeep] at (70,15) {＝};
  \node[font=\sffamily\bfseries\footnotesize, text=bkDeep] at (70,8) {同じ形};
  % 右：電流源につないだ抵抗とコンデンサー（並列。下の線が大地）
  \begin{scope}[shift={(80,0)}]
    \draw[bkInk, line width=0.9pt] (6,4) -- (6,11) (6,19) -- (6,26) -- (50,26) -- (50,18.6) (50,16.4) -- (50,4) -- (6,4);
    \draw[bkInk, line width=0.9pt] (6,15) circle (4);
    \draw[-{Stealth[length=1.8mm]}, bkInk, line width=0.8pt] (6,12.2) -- (6,17.8);
    \node[anchor=west, font=\sffamily\footnotesize] at (10.4,15) {$i$};
    \draw[bkInk, line width=0.9pt] (28,26) -- (28,21) (28,9) -- (28,4);
    \draw[bkInk, line width=0.9pt] (25.6,9) rectangle (30.4,21);
    \node[anchor=west, font=\sffamily\footnotesize] at (31,15) {$R$};
    \draw[bkInk, line width=1.1pt] (46.6,18.6) -- (53.4,18.6) (46.6,16.4) -- (53.4,16.4);
    \node[anchor=west, font=\sffamily\footnotesize] at (54,17.5) {$C$};
    \node[anchor=south, font=\sffamily\footnotesize] at (28,26.6) {電圧 $v$};
    \draw[bkInk, line width=0.8pt] (24,4) -- (32,4) (25.5,2.6) -- (30.5,2.6) (27,1.2) -- (29,1.2);
  \end{scope}
  \node[anchor=north, align=center, font=\small] at (108,-1) {$C\dfrac{\dd v}{\dd t}+\dfrac{v}{R}=i$\\[1mm]{\footnotesize $v$（大地との電位差）}};
\end{tikzpicture}
\end{bkfigure}
\vspace{1mm}
{\footnotesize\noindent
\begin{tblr}{colspec={Q[l,m,wd=19mm] X[l,m] X[l,m] X[l,m] X[l,m]}, width=\linewidth,
  row{1}={bg=bkVol!14, font=\sffamily\bfseries\footnotesize}, hlines={0.4pt,bkLine}, rowsep=1pt, colsep=3pt,
  column{1}={font=\sffamily\bfseries}}
役わり & 力学（直線の動き） & 電気 & 熱 & 流体（管の流れ） \\
流れるもの & 速度 $v$ & 電流 $i$ & 熱流 $\dot Q$ & 体積流量 $Q$ \\
押すもの & 力 $F$ & 電圧 $e$ & 温度差 $\Delta T$ & 圧力差 $\Delta p$ \\
ためるもの & ばねのやわらかさ $1/k$ & 静電容量 $C$ & 熱容量 $C_{\mathrm{th}}$ & タンク $A/(\rho g)$ \\
慣性 & 質量 $m$ & インダクタンス $L$ & （ない） & 管の中の流体 $\rho l/A$ \\
失うもの & ダンパー $c$ & 抵抗 $R$ & 熱抵抗 $R_{\mathrm{th}}$ & 管路の抵抗 \\
\end{tblr}\par}
\vspace{1.5mm}
{\footnotesize\noindent\textbf{出会う章}\quad \ph@cp{熱-1}\quad \ph@cp{熱-11}\quad \ph@cp{熱-12}（冷める）\quad \ph@cp{波-4}（温度の制御）\quad \ph@cp{電-9}（回路）\quad 解き方（6-4）\par}
\begin{seikaku}
\begin{itemize}[leftmargin=1.2\zw, itemsep=1pt, topsep=1pt, label={\color{bkNavy}●}]
\item 熱と流体の列は「差の量と電圧」、力学の列は「力と電圧」が組です。だから熱容量とタンクは、お湯の冷め方と回路をくらべた図のように、大地につないだコンデンサー、質量・ばね・ダンパーは直列の回路になります。
\item 熱流はすでに1秒あたりのエネルギーなので、温度差×熱流は仕事率になりません。熱には「慣性」がないので、熱だけの回路では、ばねとおもりのようなゆれは起こりません。
\item ふく射（熱-13）や沸騰（熱-12）、乱流の管路（力-14）では、流れが差の量に比例しないので、直線の関係は「その近く」だけで使えます。
\item 熱-12 の「相似則」（模型と本物の流れを同じ形にする条件、力-15）は、ここでの相似とは別の意味です。
\end{itemize}
\end{seikaku}
\clearpage\global\bk@decotrue}

%% ============================================================
%% p.6〜7  この本の使い方（物理の本）
%% ============================================================
\newcommand\ph@howstep[7]{% x, y, アイコン, 色, 名前, 説明, ノード名
  \node[draw=#4, line width=1pt, fill=white, rounded corners=2.5mm,
    minimum width=30.5mm, minimum height=33mm, anchor=north] (#7) at (#1,#2) {};
  \node[anchor=north] (#7i) at ([yshift=-2.2mm]#7.north) {\bkicon[1.7em]{#3}{#4}};
  \node[anchor=north, font=\sffamily\bfseries\small, text=#4, inner sep=0pt, align=center,
    minimum height=8.5mm, text width=29.5mm] (#7l) at ([yshift=-1mm]#7i.south) {#5};
  \node[anchor=north, font=\sffamily\footnotesize, text width=28mm, align=center, text=bkInk, inner sep=0pt]
    at ([yshift=-1.0mm]#7l.south) {#6};}
% 新しい囲みの見本の札  \ph@newbox{アイコン}{色}{名前}{説明}
\newcommand\ph@newbox[4]{%
  \begin{tikzpicture}[x=1mm, y=1mm]
    \draw[#2, line width=0.9pt, fill=white, rounded corners=2mm] (0,0) rectangle (67,14);
    \node[anchor=west, inner sep=0pt] (i) at (2.5,7) {\bkicon[1.5em]{#1}{#2}};
    \node[anchor=north west, inner sep=0pt, text=#2!75!black, font=\sffamily\bfseries\small] at (11.5,12) {#3};
    \node[anchor=north west, inner sep=0pt, text width=54mm, align=left, text=bkInk, font=\sffamily\fontsize{8.5}{12}\selectfont] at (11.5,6.6) {#4};
  \end{tikzpicture}}
\newcommand\phHowToPages{%
\clearpage\setcounter{page}{6}%
\global\bk@decofalse\bkfolio{true}\bksetrunning{}{}{}%
{\pdfbookmark[0]{この本の使い方}{ph-front-howto}\typeout{BKPAGE front この本の使い方 \thepage}}%
\ph@fronthead{この本の使い方}{}
{\fontsize{10.5}{19}\selectfont
この本は、身の周りの「なぜ？」に物理で答える本です。できごとを、何かを無視した「モデル」にして式を立て、答えを数で出して確かめます。式は数学シリーズの道具で書きます。\par}

\bksubsection{1}{1つの章は、いつも同じ順番}
\begin{bkfigure}
\begin{tikzpicture}[x=1mm, y=1mm]
  \ph@howstep{0}{0}{hatena}{bkYellow!85!black}{はてな？}{身近なふしぎ\\から始めよう}{s1}
  \ph@howstep{35}{0}{yattemiyou}{bkGreen}{やってみよう}{実験して、\\はかってみよう}{s2}
  \ph@howstep{70}{0}{phmodel}{phModel}{モデル化の札}{何を無視するかを\\決めよう}{s3}
  \ph@howstep{105}{0}{phflow}{bkTheme}{一般の式から、\\よく使う式へ}{条件で、よく使う\\式を導く}{s4}
  \ph@howstep{0}{-38}{reidai}{bkBlue}{例題}{7つの段で、\\ていねいに解く}{s5}
  \ph@howstep{35}{-38}{tsumazuki}{bkRed}{まちがえやすい\\ところ}{理由を知ろう}{s6}
  \ph@howstep{70}{-38}{phgear}{bkPurple}{工学ではこう使う}{本物の機械の\\数で確かめよう}{s7}
  \ph@howstep{105}{-38}{check}{bkTheme}{章末問題}{基本とチャレンジ\\の2段だよ}{s8}
  \foreach \a/\b in {s1/s2,s2/s3,s3/s4,s5/s6,s6/s7,s7/s8}
    {\draw[bk arrow, bkGray, shorten <=1mm, shorten >=1mm] (\a.east) -- (\b.west);}
  \draw[bk arrow, bkGray, rounded corners=2mm, shorten <=1.5mm, shorten >=1.5mm] (s4.south) -- ++(0,-2.5) -| (s5.north);
\end{tikzpicture}
\end{bkfigure}
\vspace{1mm}
例題は「場面→何を無視するか→図と記号→使う式を選ぶ理由→式を立てる→解く→答えを確かめる」の7つの段で解きます。こまったら、例題にもどろう。

\bksubsection{2}{章のはじめの「なぜ？」と「道しるべ」}
\leavevmode 章の帯には、その章で答える「なぜ？」があります。帯のすぐ下には、こんな行があります。
\par\vspace{1.5mm}
\noindent\begin{minipage}{\linewidth}\bkmichishirube{3-7,5-14,力-12}{熱-4,波-9}\end{minipage}
数字の札（3-7）は数学シリーズ、漢字の札（力-12・熱-4）は物理シリーズの章です。わからない式は、数学の本や第1巻『力学』にもどって確かめよう。

\clearpage
\bksubsection{3}{式が出てきたら、3つの見方で読む}
\begin{phyomi}[$Q=mc\,\Delta T$]
熱 $Q$ は、質量 $m$ と上がる温度 $\Delta T$ に比例。比例の係数が比熱 $c$。\par\noindent\phdim{J ＝ kg × J/(kg·K) × K}\enspace{\small 左と右の単位がそろう}
\end{phyomi}
\vspace{-2mm}
\begin{phmijika}
水\qty{1}{L}を20\,℃から100\,℃にわかすと約\qty{340}{kJ}。\qty{1250}{W}で約4.5分。
\end{phmijika}
\vspace{-2mm}
\begin{phbaisuru}{質量 $m$}
要る熱も、同じヒーターでわかす時間も2倍。$\Delta T$ を2倍にしても熱は2倍。
\end{phbaisuru}

\bksubsection{4}{物理の本で足した囲み}
\noindent\ph@newbox{phmodel}{phModel}{モデル化の札}{無視したこと・仮定・くずれるとき}\hfill
\ph@newbox{phketa}{phKeta}{桁の感覚}{本物の値と、答えの桁の確かめ方}\par\vspace{1mm}
\noindent\ph@newbox{phstation}{phInk}{工学の駅}{この先の工学の科目と専門書のことば}\hfill
\ph@newbox{phflow}{bkTheme}{一般の式から、よく使う式へ}{条件をおいて、よく使う式を導く}\par\vspace{1mm}
\noindent\ph@newbox{phmada}{phMada}{まだわかっていないこと}{今の科学でも答えが出ていないこと}\hfill
\ph@newbox{phsaki}{phSaki}{この先の学び}{この本の範囲をこえる理論の入口}\par

\bksubsection{5}{★と3つの段}
★は内容の高さ（★中学理科、★★高校物理基礎、★★★高校物理、★★★★大学の初め、★★★★★大学の専門）。もくじの★☆は、章のはじめ（★）から、いちばん高い所（☆）まで。3段の章は\textbf{目でわかる}→\textbf{手でわかる}→\textbf{式でわかる}で、「ここまでで十分」の旗まで読めばゴールです。

\bksubsection{6}{この本の3つのやくそく}
\begin{enumerate}[label={\bkcircled[bkNavy]{\arabic*}}, leftmargin=2.4\zw, itemsep=2pt, topsep=2pt]
\item 式を立てる前に、図をかいて、何を無視するかを書こう。
\item 答えには単位をつけ（式の中の温度は、いつも K）、本物の値と桁が合っているかを確かめよう。
\item 繰り返しの計算は、表計算やプログラムにまかせよう（シミュレーション）。
\end{enumerate}
\clearpage\global\bk@decotrue}
\makeatother
